% answers_short.tex —— 2006 年·普通高等学校招生考试·天津卷理科 · 简短答案速查
% 由 tools/make_short_answers.py 自 answers.tex 的 \gksolhead/\gkanswer 提取，
% 并与 2006/2006_数学真题_天津理_答案要点.md 逐题比对后生成，勿手改（改源 MD 后重跑）。

\begin{gknotice}
\textbf{使用说明}：本节只给最终答案便于快速核对，完整推导见解析版。答案为编写者独立推导，\textbf{非官方标准答案}。
\end{gknotice}

\begin{gknotice}
选择题答案中的 U 表示待核，未给出猜测答案，详情见解析版。
\end{gknotice}

\gkpart{选择题}
\begin{center}\begin{tabular}{ccccccccccc}
\toprule
题号 & 1 & 2 & 3 & 4 & 5 & 6 & 7 & 8 & 9 & 10 \\
\midrule
答案 & A & U & B & B & A & B & C & D & A & D \\
\bottomrule
\end{tabular}\end{center}

\gkpart{填空题}
\begin{itemize}[leftmargin=2.6em]
  \item[\textbf{11.}] 280。
  \item[\textbf{12.}] \(3/\sqrt{10}\)。
  \item[\textbf{13.}] \(3/4\)。
  \item[\textbf{14.}] 0。
  \item[\textbf{15.}] 20 吨。
  \item[\textbf{16.}] 1。
\end{itemize}

\gkpart{解答题}
\begin{itemize}[leftmargin=2.6em]
  \item[\textbf{17.}] (1) \(AB=\sqrt2\)；(2) \(\sin(2A+C)=3\sqrt7/8\)
  \item[\textbf{18.}] (1) \(63/125\)；(2) \(162/625\)；(3) \(P(\xi=k)=\binom{k-1}{2}(3/5)^3(2/5)^{k-3}\ (k=3,4,\ldots)\)
  \item[\textbf{19.}] (1) 见解析版；(2) 见解析版
  \item[\textbf{20.}] (1) 无极值；(2) \(\theta\in(\pi/6,\pi/2)\cup(3\pi/2,11\pi/6)\)；(3) \(a\in(-\infty,0]\cup[1/2+\sqrt3/8,1)\)
  \item[\textbf{21.}] (1) \(\lambda=\pm1\)；(2) 见解析版；(3) 见解析版
  \item[\textbf{22.}] (1) \(c^2=ab\)，且 \(M=(0,a)\)；(2) 见解析版
\end{itemize}
